\documentclass[11pt,a4paper]{article}

% 한글 지원
\usepackage{fontspec}
\usepackage{xeCJK}
\setCJKmainfont{Malgun Gothic}

% 수학 패키지
\usepackage{amsmath}
\usepackage{amssymb}
\usepackage{amsthm}

% 레이아웃
\usepackage[margin=2.5cm]{geometry}
\usepackage{enumitem}

% 색상
\usepackage{xcolor}

% 표
\usepackage{booktabs}
\usepackage{array}

% 하이퍼링크
\usepackage{hyperref}
\hypersetup{
    colorlinks=true,
    linkcolor=blue,
    filecolor=magenta,
    urlcolor=cyan,
}

% 정리 환경
\newtheorem{theorem}{정리}[section]
\newtheorem{lemma}[theorem]{보조 정리}
\newtheorem{proposition}[theorem]{명제}
\theoremstyle{definition}
\newtheorem{definition}[theorem]{정의}
\newtheorem{example}[theorem]{예제}
\newtheorem{exercise}{연습문제}[section]
\theoremstyle{remark}
\newtheorem*{solution}{풀이}
\newtheorem*{remark}{주의}

% 제목 설정
\title{\textbf{고급 정수론 기법:\\카마이클 함수, p-adic 분석,\\중국인의 나머지 정리}}
\author{}
\date{2025-11-17}

\begin{document}

\maketitle

\tableofcontents
\newpage

\section{카마이클 함수 (Carmichael's Lambda Function)}

\subsection{정의와 기본 개념}

\begin{definition}[카마이클 함수]
\textbf{카마이클 함수} $\lambda(n)$은 다음을 만족하는 \textbf{가장 작은} 양의 정수입니다:
\[a^{\lambda(n)} \equiv 1 \pmod{n} \quad \text{for all } \gcd(a,n) = 1\]

이는 \textbf{군론}에서 말하는 $(\mathbb{Z}/n\mathbb{Z})^*$의 \textbf{지수(exponent)}입니다.
\end{definition}

\subsubsection{오일러 함수와의 관계}

\textbf{오일러 정리}는 다음을 말합니다:
\[a^{\phi(n)} \equiv 1 \pmod{n} \quad \text{for } \gcd(a,n) = 1\]

카마이클 함수는 이것의 \textbf{개선된 버전}입니다:
\[\lambda(n) \mid \phi(n) \quad \text{and} \quad \lambda(n) \leq \phi(n)\]

즉, 카마이클 함수는 1이 되는 \textbf{최소} 지수를 제공합니다.

\subsection{계산 공식}

\subsubsection{규칙 1: 소수 거듭제곱}

\textbf{2의 거듭제곱 (특수 케이스)}:
\begin{itemize}
    \item $\lambda(1) = 1$
    \item $\lambda(2) = 1$
    \item $\lambda(4) = 2$
    \item $\lambda(2^k) = 2^{k-2}$ for $k \geq 3$
\end{itemize}

\textbf{홀수 소수 $p$의 거듭제곱}:
\[\lambda(p^k) = \phi(p^k) = p^{k-1}(p-1)\]

\subsubsection{규칙 2: 서로소인 수들의 곱}

$\gcd(m,n) = 1$일 때:
\[\lambda(mn) = \text{lcm}(\lambda(m), \lambda(n))\]

\begin{remark}
오일러 함수는 곱셈적이지만 ($\phi(mn) = \phi(m)\phi(n)$), 카마이클 함수는 \textbf{LCM}을 사용합니다!
\end{remark}

\subsection{예시: $\lambda(100)$ 계산}

$100 = 4 \times 25 = 2^2 \times 5^2$

\textbf{단계 1}: 각 소수 거듭제곱에 대해:
\begin{align*}
\lambda(4) &= \lambda(2^2) = 2 \\
\lambda(25) &= \lambda(5^2) = \phi(5^2) = 5^{2-1} \cdot (5-1) = 5 \cdot 4 = 20
\end{align*}

\textbf{단계 2}: LCM 계산:
\[\lambda(100) = \text{lcm}(2, 20) = 20\]

\subsubsection{오일러 함수와의 비교}

\[\phi(100) = \phi(4) \cdot \phi(25) = 2 \cdot 20 = 40\]

\textbf{결과}:
\begin{itemize}
    \item $\lambda(100) = 20$ $\leftarrow$ 더 작음!
    \item $\phi(100) = 40$
\end{itemize}

이것은 $\gcd(a, 100) = 1$인 모든 $a$에 대해:
\[a^{20} \equiv 1 \pmod{100}\]
\[a^{40} \equiv 1 \pmod{100}\]

둘 다 맞지만, \textbf{20이 최소 지수}입니다.

\subsection{검증 예시}

$3$과 $100$은 서로소입니다 ($\gcd(3, 100) = 1$).

\textbf{카마이클 함수로}:
\[3^{20} \bmod 100 = ?\]

단계별 계산:
\begin{align*}
3^1 &= 3 \\
3^2 &= 9 \\
3^4 &= 81 \\
3^5 &= 243 \equiv 43 \pmod{100} \\
3^{10} &= (3^5)^2 \equiv 43^2 = 1849 \equiv 49 \pmod{100} \\
3^{20} &= (3^{10})^2 \equiv 49^2 = 2401 \equiv 1 \pmod{100} \quad \checkmark
\end{align*}

\textbf{오일러 함수로도 확인}:
\[3^{40} = (3^{20})^2 \equiv 1^2 = 1 \pmod{100} \quad \checkmark\]

\subsection{더 많은 예시}

\begin{table}[h]
\centering
\begin{tabular}{cccc}
\toprule
$n$ & 소인수분해 & $\lambda(n)$ & $\phi(n)$ \\
\midrule
12 & $2^2 \times 3$ & lcm(2, 2) = 2 & 4 \\
15 & $3 \times 5$ & lcm(2, 4) = 4 & 8 \\
20 & $2^2 \times 5$ & lcm(2, 4) = 4 & 8 \\
21 & $3 \times 7$ & lcm(2, 6) = 6 & 12 \\
24 & $2^3 \times 3$ & lcm(2, 2) = 2 & 8 \\
30 & $2 \times 3 \times 5$ & lcm(1, 2, 4) = 4 & 8 \\
\bottomrule
\end{tabular}
\caption{카마이클 함수와 오일러 함수 비교}
\end{table}

\textbf{패턴 관찰}:
\begin{itemize}
    \item 2의 거듭제곱 때문에 $\lambda(n)$이 $\phi(n)$보다 훨씬 작을 수 있습니다.
    \item $n$이 홀수이고 제곱 없는 수(square-free)이면 $\lambda(n) = \phi(n)/2$
\end{itemize}

\newpage

\section{p-adic 분석 (p-adic Valuation)}

\subsection{정의와 기본 개념}

\begin{definition}[p-adic valuation]
\textbf{$p$-adic valuation} $v_p(n)$은 정수 $n$을 나누는 소수 $p$의 \textbf{최대 거듭제곱 지수}입니다.

수학적으로:
\[n = p^{v_p(n)} \cdot m \quad \text{where } \gcd(p, m) = 1\]
\end{definition}

\subsubsection{표기법}

\begin{itemize}
    \item $v_p(n)$: $p$-adic valuation
    \item $\nu_p(n)$: 다른 표기법
    \item $\text{ord}_p(n)$: 또 다른 표기법
\end{itemize}

\begin{example}[p-adic valuation 계산]
\begin{enumerate}
    \item $v_2(24) = ?$
    \begin{itemize}
        \item $24 = 2^3 \times 3$
        \item $v_2(24) = 3$
    \end{itemize}

    \item $v_5(250) = ?$
    \begin{itemize}
        \item $250 = 2 \times 5^3$
        \item $v_5(250) = 3$
    \end{itemize}

    \item $v_3(100) = ?$
    \begin{itemize}
        \item $100 = 2^2 \times 5^2$
        \item $v_3(100) = 0$ (3으로 나누어지지 않음)
    \end{itemize}

    \item $v_7(49000) = ?$
    \begin{itemize}
        \item $49000 = 49 \times 1000 = 7^2 \times 2^3 \times 5^3$
        \item $v_7(49000) = 2$
    \end{itemize}
\end{enumerate}
\end{example}

\subsection{기본 성질}

\begin{proposition}[p-adic valuation의 성질]
\begin{enumerate}
    \item \textbf{곱셈}: $v_p(ab) = v_p(a) + v_p(b)$
    \item \textbf{거듭제곱}: $v_p(a^k) = k \cdot v_p(a)$
    \item \textbf{나눗셈}: $v_p(a/b) = v_p(a) - v_p(b)$ (단, $b \neq 0$)
    \item \textbf{덧셈 (삼각 부등식)}: $v_p(a + b) \geq \min(v_p(a), v_p(b))$
\end{enumerate}
\end{proposition}

\begin{proof}[성질 1의 증명]
\begin{itemize}
    \item $a = p^{\alpha} a'$, $\gcd(p, a') = 1$
    \item $b = p^{\beta} b'$, $\gcd(p, b') = 1$
    \item $ab = p^{\alpha + \beta} a'b'$, $\gcd(p, a'b') = 1$
    \item $\therefore v_p(ab) = \alpha + \beta = v_p(a) + v_p(b)$ \qed
\end{itemize}
\end{proof}

\begin{remark}
덧셈에서 등호 조건: $v_p(a) \neq v_p(b)$이면
\[v_p(a + b) = \min(v_p(a), v_p(b))\]
\end{remark}

\subsection{르장드르 공식 (Legendre's Formula)}

\begin{theorem}[르장드르 공식]
$N!$에서 소수 $p$의 거듭제곱:
\[v_p(N!) = \sum_{i=1}^{\infty} \left\lfloor \frac{N}{p^i} \right\rfloor\]

실제로는 $p^i > N$이면 항이 0이므로 유한 합입니다.
\end{theorem}

\begin{example}[$v_5(100!)$ 계산]
\begin{align*}
v_5(100!) &= \left\lfloor \frac{100}{5} \right\rfloor + \left\lfloor \frac{100}{25} \right\rfloor + \left\lfloor \frac{100}{125} \right\rfloor + \cdots \\
&= 20 + 4 + 0 + \cdots = 24
\end{align*}

따라서 $100! = 5^{24} \times (\text{5와 서로소인 수})$
\end{example}

\subsubsection{직관적 이해}

\begin{itemize}
    \item $\lfloor 100/5 \rfloor = 20$: 5의 배수 개수
    \item $\lfloor 100/25 \rfloor = 4$: 25의 배수 개수 (추가로 5를 하나 더)
    \item $\lfloor 100/125 \rfloor = 0$: 125의 배수 개수
\end{itemize}

각 $p^i$의 배수는 $p$를 $i$번 포함하지만, 이미 이전 단계에서 센 것을 빼야 합니다.
합으로 표현하면 자동으로 처리됩니다.

\subsection{원래 문제에 적용}

문제: $P = 1^2 \cdot 3^2 \cdot 5^2 \cdot 7^2 \cdots 99999^2$

$P \bmod 100$을 알려면 $v_2(P)$와 $v_5(P)$를 알아야 합니다.

\subsubsection{$v_2(P)$ 계산}

모든 홀수는 2로 나누어지지 않으므로:
\[v_2(\text{홀수}) = 0\]
\[v_2(P) = 0\]

\subsubsection{$v_5(P)$ 계산}

\textbf{단계 1}: 1부터 99999까지 홀수 중 5의 배수 찾기

홀수는 $1, 3, 5, 7, \ldots, 99999$ (총 50000개)

5의 배수인 홀수: $5, 15, 25, 35, \ldots, 99995$

이들은 $5(2k-1)$ 형태이고, $k = 1, 2, 3, \ldots, 10000$

따라서 \textbf{10000개}입니다.

\textbf{단계 2}: 홀수들의 곱에서 $v_5$ 계산

홀수들의 곱 $Q = \prod_{k=1}^{50000} (2k-1)$

\textbf{르장드르 공식 수정 버전}:

홀수들의 곱에서 5의 거듭제곱:
\begin{align*}
v_5(Q) &= \left\lfloor \frac{50000}{5} \right\rfloor + \left\lfloor \frac{50000}{25} \right\rfloor + \left\lfloor \frac{50000}{125} \right\rfloor + \left\lfloor \frac{50000}{625} \right\rfloor \\
&\quad + \left\lfloor \frac{50000}{3125} \right\rfloor + \left\lfloor \frac{50000}{15625} \right\rfloor \\
&= 10000 + 2000 + 400 + 80 + 16 + 3 \\
&= 12499
\end{align*}

\textbf{단계 3}: 제곱 고려
\[P = Q^2\]
\[v_5(P) = 2 \times v_5(Q) = 2 \times 12499 = 24998\]

\subsubsection{결론}

\[P = 2^0 \times 5^{24998} \times (\text{2, 5와 서로소인 수})\]

따라서:
\begin{itemize}
    \item $v_5(P) = 24998 \geq 2$
    \item $P \equiv 0 \pmod{25}$
    \item $P \equiv 1 \pmod{4}$ (홀수 제곱들의 곱)
\end{itemize}

\textbf{중국인의 나머지 정리}:
\[\begin{cases}
P \equiv 1 \pmod{4} \\
P \equiv 0 \pmod{25}
\end{cases}\]

$P = 25k$이고 $25k \equiv 1 \pmod{4}$

$k \equiv 1 \pmod{4}$

$k = 4m + 1$

$P = 25(4m + 1) = 100m + 25$

\[\boxed{P \equiv 25 \pmod{100}}\]

\newpage

\section{중국인의 나머지 정리 (Chinese Remainder Theorem)}

\subsection{정의와 정리}

\begin{theorem}[중국인의 나머지 정리]
$m_1, m_2, \ldots, m_k$가 쌍으로 서로소인 양의 정수들이고 (즉, $\gcd(m_i, m_j) = 1$ for $i \neq j$),
$a_1, a_2, \ldots, a_k$가 임의의 정수들일 때, 다음 연립 합동식:

\[\begin{cases}
x \equiv a_1 \pmod{m_1} \\
x \equiv a_2 \pmod{m_2} \\
\vdots \\
x \equiv a_k \pmod{m_k}
\end{cases}\]

은 $\bmod M = m_1 \cdot m_2 \cdot \cdots \cdot m_k$에서 \textbf{유일한 해}를 가집니다.
\end{theorem}

\subsubsection{역사적 배경}

이 정리는 중국 수학자 \textbf{손자(孫子, Sun Tzu, 4세기경)}의 저서 『손자산경(孫子算經)』에 처음 등장했습니다.

\textbf{원래 문제}:
``물건의 개수를 3으로 나누면 2가 남고, 5로 나누면 3이 남고, 7로 나누면 2가 남는다. 물건은 몇 개인가?''

\[\begin{cases}
x \equiv 2 \pmod{3} \\
x \equiv 3 \pmod{5} \\
x \equiv 2 \pmod{7}
\end{cases}\]

답: $x \equiv 23 \pmod{105}$ (최소 양수 해는 23)

\subsection{두 개 합동식의 경우}

\subsubsection{기본 사례: 2개 연립 합동식}

\[\begin{cases}
x \equiv a_1 \pmod{m_1} \\
x \equiv a_2 \pmod{m_2}
\end{cases}\]

여기서 $\gcd(m_1, m_2) = 1$

\subsubsection{구성적 증명 (알고리즘)}

\textbf{단계 1}: 각 모듈러스에 대한 ``역원'' 찾기

$M = m_1 \cdot m_2$라 하고:
\begin{itemize}
    \item $M_1 = M/m_1 = m_2$
    \item $M_2 = M/m_2 = m_1$
\end{itemize}

\textbf{단계 2}: Bézout 항등식 사용

$\gcd(M_1, m_1) = \gcd(m_2, m_1) = 1$이므로, 다음을 만족하는 $y_1$이 존재:
\[M_1 \cdot y_1 \equiv 1 \pmod{m_1}\]
즉, $m_2 \cdot y_1 \equiv 1 \pmod{m_1}$

마찬가지로:
\[M_2 \cdot y_2 \equiv 1 \pmod{m_2}\]
즉, $m_1 \cdot y_2 \equiv 1 \pmod{m_2}$

\textbf{단계 3}: 해 구성

\[x = a_1 \cdot M_1 \cdot y_1 + a_2 \cdot M_2 \cdot y_2\]

\begin{example}[간단한 2개 합동식]
\[\begin{cases}
x \equiv 2 \pmod{5} \\
x \equiv 3 \pmod{7}
\end{cases}\]

$M = 5 \times 7 = 35$

$M_1 = 35/5 = 7$, $M_2 = 35/7 = 5$

\textbf{$y_1$ 찾기}: $7 \cdot y_1 \equiv 1 \pmod{5}$
\begin{itemize}
    \item $7 \equiv 2 \pmod{5}$
    \item $2 \cdot y_1 \equiv 1 \pmod{5}$
    \item $y_1 = 3$ (왜냐하면 $2 \times 3 = 6 \equiv 1 \pmod{5}$)
\end{itemize}

\textbf{$y_2$ 찾기}: $5 \cdot y_2 \equiv 1 \pmod{7}$
\begin{itemize}
    \item $y_2 = 3$ (왜냐하면 $5 \times 3 = 15 \equiv 1 \pmod{7}$)
\end{itemize}

\textbf{해 구성}:
\[x = 2 \cdot 7 \cdot 3 + 3 \cdot 5 \cdot 3 = 42 + 45 = 87\]
\[87 \bmod 35 = 17\]

\textbf{답}: $x \equiv 17 \pmod{35}$

\textbf{검증}:
\begin{itemize}
    \item $17 = 3 \times 5 + 2 \equiv 2 \pmod{5}$ \checkmark
    \item $17 = 2 \times 7 + 3 \equiv 3 \pmod{7}$ \checkmark
\end{itemize}
\end{example}

\subsection{세 개 이상 합동식의 경우}

\subsubsection{일반 알고리즘}

\[\begin{cases}
x \equiv a_1 \pmod{m_1} \\
x \equiv a_2 \pmod{m_2} \\
x \equiv a_3 \pmod{m_3}
\end{cases}\]

\textbf{방법 1}: 순차적 해결
\begin{enumerate}
    \item 먼저 처음 두 개를 풀어 $x \equiv c \pmod{m_1 \cdot m_2}$를 얻음
    \item 이것과 세 번째 식을 결합하여 최종 해를 구함
\end{enumerate}

\textbf{방법 2}: 직접 구성

$M = m_1 \cdot m_2 \cdot m_3$

각 $i$에 대해:
\begin{itemize}
    \item $M_i = M / m_i$
    \item $M_i \cdot y_i \equiv 1 \pmod{m_i}$를 만족하는 $y_i$ 찾기
\end{itemize}

해:
\[x = \sum_{i=1}^{3} a_i \cdot M_i \cdot y_i\]

\begin{example}[손자의 원래 문제]
\[\begin{cases}
x \equiv 2 \pmod{3} \\
x \equiv 3 \pmod{5} \\
x \equiv 2 \pmod{7}
\end{cases}\]

$M = 3 \times 5 \times 7 = 105$

$M_1 = 105/3 = 35$, $M_2 = 105/5 = 21$, $M_3 = 105/7 = 15$

\textbf{$y_1$ 찾기}: $35 \cdot y_1 \equiv 1 \pmod{3}$
\begin{itemize}
    \item $35 = 11 \times 3 + 2 \equiv 2 \pmod{3}$
    \item $2 \cdot y_1 \equiv 1 \pmod{3}$
    \item $y_1 = 2$ (왜냐하면 $2 \times 2 = 4 \equiv 1 \pmod{3}$)
\end{itemize}

\textbf{$y_2$ 찾기}: $21 \cdot y_2 \equiv 1 \pmod{5}$
\begin{itemize}
    \item $21 = 4 \times 5 + 1 \equiv 1 \pmod{5}$
    \item $y_2 = 1$
\end{itemize}

\textbf{$y_3$ 찾기}: $15 \cdot y_3 \equiv 1 \pmod{7}$
\begin{itemize}
    \item $15 = 2 \times 7 + 1 \equiv 1 \pmod{7}$
    \item $y_3 = 1$
\end{itemize}

\textbf{해 구성}:
\begin{align*}
x &= 2 \cdot 35 \cdot 2 + 3 \cdot 21 \cdot 1 + 2 \cdot 15 \cdot 1 \\
  &= 140 + 63 + 30 = 233
\end{align*}

$233 = 2 \times 105 + 23$

\textbf{답}: $x \equiv 23 \pmod{105}$

\textbf{검증}:
\begin{itemize}
    \item $23 = 7 \times 3 + 2 \equiv 2 \pmod{3}$ \checkmark
    \item $23 = 4 \times 5 + 3 \equiv 3 \pmod{5}$ \checkmark
    \item $23 = 3 \times 7 + 2 \equiv 2 \pmod{7}$ \checkmark
\end{itemize}

최소 양의 정수 해는 \textbf{23}입니다!
\end{example}

\subsection{응용과 활용}

\begin{enumerate}
    \item \textbf{큰 수의 모듈러 연산}: 큰 모듈러스 $n = p_1 \cdot p_2 \cdot \cdots \cdot p_k$ (소인수분해)에 대한 계산을 각 소인수별로 분리하여 계산 후 결합
    \item \textbf{달력 문제}: 여러 주기가 있는 현상들의 동시 발생 시점 찾기
    \item \textbf{수 이론 문제}: 여러 나머지 조건을 만족하는 수 찾기
    \item \textbf{암호학}: RSA 암호 시스템의 핵심 도구
\end{enumerate}

\newpage

\section{연습문제}

\subsection{카마이클 함수 연습문제}

\begin{exercise}[$\lambda(360)$ 계산]
$\lambda(360)$을 계산하시오.
\end{exercise}

\begin{solution}

\textbf{단계 1}: 소인수분해
\[360 = 8 \times 45 = 2^3 \times 3^2 \times 5\]

\textbf{단계 2}: 각 소수 거듭제곱에 대한 카마이클 함수
\begin{align*}
\lambda(2^3) &= 2^{3-2} = 2 \\
\lambda(3^2) &= \phi(3^2) = 3^{2-1}(3-1) = 3 \times 2 = 6 \\
\lambda(5) &= \phi(5) = 4
\end{align*}

\textbf{단계 3}: LCM 계산
\[\lambda(360) = \text{lcm}(2, 6, 4)\]

$2 = 2$, $6 = 2 \times 3$, $4 = 2^2$

\[\text{lcm}(2, 6, 4) = 2^2 \times 3 = 12\]

\textbf{답}: $\lambda(360) = 12$

\textbf{검증}: $\phi(360) = \phi(2^3) \times \phi(3^2) \times \phi(5) = 4 \times 6 \times 4 = 96$

$\lambda(360) = 12$는 $\phi(360) = 96$의 약수입니다. ($96 = 12 \times 8$) \checkmark

\textbf{의미}: $\gcd(a, 360) = 1$인 모든 $a$에 대해:
\[a^{12} \equiv 1 \pmod{360}\]
\end{solution}

\begin{exercise}[최소 지수 찾기]
$7^k \equiv 1 \pmod{100}$을 만족하는 \textbf{최소 양의 정수} $k$를 구하시오.
\end{exercise}

\begin{solution}

\textbf{단계 1}: $\gcd(7, 100) = 1$ 확인 \checkmark

\textbf{단계 2}: $\lambda(100)$ 계산

앞에서 계산했듯이:
\[\lambda(100) = \lambda(4 \times 25) = \text{lcm}(\lambda(4), \lambda(25)) = \text{lcm}(2, 20) = 20\]

따라서 $7^{20} \equiv 1 \pmod{100}$

\textbf{단계 3}: 최소 지수는 $\lambda(100)$의 약수

$k$는 20의 약수 중 하나: $1, 2, 4, 5, 10, 20$

\textbf{단계 4}: 각각 확인

\begin{align*}
7^1 &= 7 \not\equiv 1 \pmod{100} \\
7^2 &= 49 \not\equiv 1 \pmod{100} \\
7^4 &= 2401 \equiv 1 \pmod{100} \quad \checkmark
\end{align*}

\textbf{답}: $k = 4$

\textbf{상세 계산}:
\[7^4 = 49^2 = 2401 = 24 \times 100 + 1 \equiv 1 \pmod{100}\] \checkmark

\textbf{추가 관찰}:
\begin{itemize}
    \item $7^4 \equiv 1 \pmod{100}$이므로 7의 위수(order)는 4입니다.
    \item $\lambda(100) = 20 = 4 \times 5$이므로 일관성이 있습니다.
    \item 7의 거듭제곱 mod 100의 주기는 4입니다.
\end{itemize}
\end{solution}

\subsection{p-adic 분석 연습문제}

\begin{exercise}[팩토리얼의 끝자리 0의 개수]
$1000!$의 끝자리에 0이 몇 개 붙는가?
\end{exercise}

\begin{solution}

끝자리 0의 개수 = $\min(v_2(1000!), v_5(1000!))$

일반적으로 $v_5(n!) < v_2(n!)$이므로 $v_5(1000!)$만 계산하면 됩니다.

\textbf{르장드르 공식}:
\[v_5(1000!) = \sum_{i=1}^{\infty} \left\lfloor \frac{1000}{5^i} \right\rfloor\]

계산:
\begin{align*}
\left\lfloor \frac{1000}{5} \right\rfloor &= 200 \\
\left\lfloor \frac{1000}{25} \right\rfloor &= 40 \\
\left\lfloor \frac{1000}{125} \right\rfloor &= 8 \\
\left\lfloor \frac{1000}{625} \right\rfloor &= 1 \\
\left\lfloor \frac{1000}{3125} \right\rfloor &= 0
\end{align*}

합계:
\[v_5(1000!) = 200 + 40 + 8 + 1 = 249\]

\textbf{답}: $1000!$의 끝자리에는 \textbf{249개의 0}이 붙습니다.

\textbf{검증}: $v_2(1000!)$ 계산 (참고)
\[v_2(1000!) = 500 + 250 + 125 + 62 + 31 + 15 + 7 + 3 + 1 = 994\]

$994 > 249$이므로 끝자리 0의 개수는 249개가 맞습니다. \checkmark
\end{solution}

\begin{exercise}[이항계수의 나누어떨어짐]
$\binom{2025}{1000}$이 3으로 나누어떨어지는가?
\end{exercise}

\begin{solution}

$\displaystyle\binom{2025}{1000} = \frac{2025!}{1000! \times 1025!}$

\textbf{$v_3$ 계산}:
\[v_3\left(\binom{2025}{1000}\right) = v_3(2025!) - v_3(1000!) - v_3(1025!)\]

\textbf{$v_3(2025!)$ 계산}:
\begin{align*}
v_3(2025!) &= \left\lfloor \frac{2025}{3} \right\rfloor + \left\lfloor \frac{2025}{9} \right\rfloor + \left\lfloor \frac{2025}{27} \right\rfloor + \left\lfloor \frac{2025}{81} \right\rfloor \\
&\quad + \left\lfloor \frac{2025}{243} \right\rfloor + \left\lfloor \frac{2025}{729} \right\rfloor + \left\lfloor \frac{2025}{2187} \right\rfloor \\
&= 675 + 225 + 75 + 25 + 8 + 2 + 0 = 1010
\end{align*}

\textbf{$v_3(1000!)$ 계산}:
\begin{align*}
v_3(1000!) &= 333 + 111 + 37 + 12 + 4 + 1 = 498
\end{align*}

\textbf{$v_3(1025!)$ 계산}:
\begin{align*}
v_3(1025!) &= 341 + 113 + 37 + 12 + 4 + 1 = 508
\end{align*}

\textbf{최종 계산}:
\[v_3\left(\binom{2025}{1000}\right) = 1010 - 498 - 508 = 4\]

\textbf{답}: $v_3\left(\binom{2025}{1000}\right) = 4 > 0$이므로, $\binom{2025}{1000}$은 \textbf{3으로 나누어떨어집니다}.

실제로 $3^4 = 81$로 나누어떨어집니다!
\end{solution}

\subsection{종합 문제}

\begin{exercise}[카마이클과 p-adic 결합]
$1 \times 3 \times 5 \times 7 \times \cdots \times 999$의 값을 $\bmod 1000$으로 계산하시오.
\end{exercise}

\begin{solution}

$P = 1 \times 3 \times 5 \times 7 \times \cdots \times 999$ (500개의 홀수)

\textbf{p-adic 분석}:

$1000 = 2^3 \times 5^3$

\textbf{$v_2(P)$ 계산}:
모든 홀수이므로 $v_2(P) = 0$

\textbf{$v_5(P)$ 계산}:

5의 배수인 홀수: $5, 15, 25, 35, \ldots, 995$ (250개)

르장드르 공식:
\begin{align*}
v_5(P) &= \left\lfloor \frac{500}{5} \right\rfloor + \left\lfloor \frac{500}{25} \right\rfloor + \left\lfloor \frac{500}{125} \right\rfloor + \left\lfloor \frac{500}{625} \right\rfloor \\
&= 100 + 20 + 4 + 0 = 124
\end{align*}

$v_5(P) = 124 \geq 3$이므로:
\[P \equiv 0 \pmod{125}\]
\[P \equiv 0 \pmod{1000}\]

\textbf{답}: $P \equiv 0 \pmod{1000}$

\textbf{의미}: 1부터 999까지 모든 홀수를 곱하면 마지막 \textbf{세 자리}가 모두 \textbf{000}입니다!
\end{solution}

\subsection{중국인의 나머지 정리 연습문제}

\begin{exercise}[2개 합동식 풀기]
다음 연립 합동식을 풀어라.

\[\begin{cases}
x \equiv 3 \pmod{8} \\
x \equiv 5 \pmod{13}
\end{cases}\]
\end{exercise}

\begin{solution}

$\gcd(8, 13) = 1$이므로 중국인의 나머지 정리를 적용할 수 있습니다.

$M = 8 \times 13 = 104$

$M_1 = 104/8 = 13$, $M_2 = 104/13 = 8$

\textbf{$y_1$ 찾기}: $13 \cdot y_1 \equiv 1 \pmod{8}$
\begin{itemize}
    \item $13 = 1 \times 8 + 5 \equiv 5 \pmod{8}$
    \item $5 \cdot y_1 \equiv 1 \pmod{8}$
    \item 시행착오: $5 \times 5 = 25 \equiv 1 \pmod{8}$
    \item $y_1 = 5$
\end{itemize}

\textbf{$y_2$ 찾기}: $8 \cdot y_2 \equiv 1 \pmod{13}$
\begin{itemize}
    \item 확장 유클리드: $8 \times 5 = 40 = 3 \times 13 + 1$
    \item $y_2 = 5$
\end{itemize}

\textbf{해 구성}:
\begin{align*}
x &= 3 \cdot 13 \cdot 5 + 5 \cdot 8 \cdot 5 \\
  &= 195 + 200 = 395
\end{align*}

$395 = 3 \times 104 + 83$

\textbf{답}: $x \equiv 83 \pmod{104}$

\textbf{검증}:
\begin{itemize}
    \item $83 = 10 \times 8 + 3 \equiv 3 \pmod{8}$ \checkmark
    \item $83 = 6 \times 13 + 5 \equiv 5 \pmod{13}$ \checkmark
\end{itemize}

\textbf{최소 양의 정수 해}: $x = 83$
\end{solution}

\begin{exercise}[3개 합동식 풀기]
다음 연립 합동식을 풀어라.

\[\begin{cases}
x \equiv 1 \pmod{4} \\
x \equiv 2 \pmod{5} \\
x \equiv 3 \pmod{11}
\end{cases}\]
\end{exercise}

\begin{solution}

$\gcd(4, 5) = \gcd(4, 11) = \gcd(5, 11) = 1$이므로 중국인의 나머지 정리를 적용할 수 있습니다.

$M = 4 \times 5 \times 11 = 220$

$M_1 = 220/4 = 55$, $M_2 = 220/5 = 44$, $M_3 = 220/11 = 20$

\textbf{$y_1$ 찾기}: $55 \cdot y_1 \equiv 1 \pmod{4}$
\begin{itemize}
    \item $55 = 13 \times 4 + 3 \equiv 3 \pmod{4}$
    \item $3 \cdot y_1 \equiv 1 \pmod{4}$
    \item $y_1 = 3$ (왜냐하면 $3 \times 3 = 9 \equiv 1 \pmod{4}$)
\end{itemize}

\textbf{$y_2$ 찾기}: $44 \cdot y_2 \equiv 1 \pmod{5}$
\begin{itemize}
    \item $44 = 8 \times 5 + 4 \equiv 4 \pmod{5}$
    \item $4 \cdot y_2 \equiv 1 \pmod{5}$
    \item $y_2 = 4$ (왜냐하면 $4 \times 4 = 16 \equiv 1 \pmod{5}$)
\end{itemize}

\textbf{$y_3$ 찾기}: $20 \cdot y_3 \equiv 1 \pmod{11}$
\begin{itemize}
    \item $20 = 1 \times 11 + 9 \equiv 9 \pmod{11}$
    \item $9 \cdot y_3 \equiv 1 \pmod{11}$
    \item $9 \times 5 = 45 = 4 \times 11 + 1 \equiv 1 \pmod{11}$
    \item $y_3 = 5$
\end{itemize}

\textbf{해 구성}:
\begin{align*}
x &= 1 \cdot 55 \cdot 3 + 2 \cdot 44 \cdot 4 + 3 \cdot 20 \cdot 5 \\
  &= 165 + 352 + 300 = 817
\end{align*}

$817 = 3 \times 220 + 157$

\textbf{답}: $x \equiv 157 \pmod{220}$

\textbf{검증}:
\begin{itemize}
    \item $157 = 39 \times 4 + 1 \equiv 1 \pmod{4}$ \checkmark
    \item $157 = 31 \times 5 + 2 \equiv 2 \pmod{5}$ \checkmark
    \item $157 = 14 \times 11 + 3 \equiv 3 \pmod{11}$ \checkmark
\end{itemize}

\textbf{최소 양의 정수 해}: $x = 157$
\end{solution}

\newpage

\section{핵심 요약}

\subsection{카마이클 함수}

\begin{itemize}
    \item \textbf{정의}: $a^{\lambda(n)} \equiv 1 \pmod{n}$을 만족하는 최소 지수
    \item \textbf{계산}: $\lambda(p^k)$ 공식 + LCM
    \item \textbf{활용}: 거듭제곱 계산 최적화, 주기 분석
\end{itemize}

\subsection{p-adic 분석}

\begin{itemize}
    \item \textbf{정의}: $v_p(n)$ = $n$을 나누는 $p$의 최대 거듭제곱 지수
    \item \textbf{성질}: 곱셈, 거듭제곱, 나눗셈 규칙
    \item \textbf{르장드르 공식}: $v_p(N!) = \sum_{i \geq 1} \lfloor N/p^i \rfloor$
    \item \textbf{활용}: 나누어떨어짐 분석, 조합론 문제
\end{itemize}

\subsection{중국인의 나머지 정리}

\begin{itemize}
    \item \textbf{정의}: 서로소인 모듈러스에 대한 연립 합동식의 유일해 존재
    \item \textbf{2개 합동식}: 직접 대입 또는 구성적 방법
    \item \textbf{3개 이상}: 순차적 해결 또는 일반 공식
    \item \textbf{확장 유클리드}: 모듈러 역원 찾기
    \item \textbf{활용}: 큰 수의 모듈러 연산, RSA 암호화
\end{itemize}

\subsection{결합 사용}

\begin{itemize}
    \item 세 기법을 함께 사용하여 복잡한 문제 해결
    \item 각 소수별로 독립적 분석 (p-adic) 후 결합 (CRT)
    \item 거듭제곱 최적화 (카마이클 함수)
\end{itemize}

\vspace{1cm}

\hrule

\vspace{0.5cm}

\textbf{작성일}: 2025-11-17

\textbf{주제}: 고급 정수론, 카마이클 함수, p-adic 분석, 중국인의 나머지 정리, 르장드르 공식

\end{document}
